\begin{table}[H]\centering
\caption{Contrastes exploratorios cerca del límite educativo}\label{tab:short_local}
\begin{threeparttable}\small
\setlength{\tabcolsep}{4pt}
\begin{tabular}{lrrr}
\toprule
Resultado & Coeficiente & IC 95\% & $p$ \\
\midrule
\multicolumn{4}{l}{\emph{Secundaria completa frente a superior no universitaria incompleta}} \\
Salario horario (log) & -0.025 & [-0.074, 0.024] & 0.310 \\
Formalidad (pp) & -1.403 & [-3.933, 1.127] & 0.264 \\
Autoempleo (pp) & 1.303 & [-0.271, 2.877] & 0.100 \\
\multicolumn{4}{l}{\emph{Cambio del salto lineal extrapolado en 12 años}} \\
Salario horario (log) & -0.050 & [-0.144, 0.045] & 0.288 \\
Formalidad (pp) & 1.123 & [-3.441, 5.688] & 0.616 \\
Autoempleo (pp) & -1.460 & [-4.627, 1.707] & 0.351 \\
\bottomrule
\end{tabular}
\begin{tablenotes}\footnotesize
\item Orientación baja menos alta educación. Primer panel: DiD en categorías de 11 y 13 años imputados (29\,008 observaciones salariales y 64\,143 de composición). Segundo: diferencia de saltos lineales con valores educativos 9, 11, 13 y 14 (51\,072 y 117\,450 observaciones). Pesos de encuesta, controles individuales y efectos fijos de departamento y trimestre. Errores agrupados en 25 departamentos. Estos ejercicios no identifican una RD causal; el límite no asigna la reforma.
\end{tablenotes}
\end{threeparttable}\end{table}
